\section{Cost and the compiled implementation}
\label{S-sec:s-cost}

\begin{table}[t]
\caption[Run time on 3DLoMatch+FPFH]{Median seconds per pair on an AMD EPYC 9655P, using the same 100 3DLoMatch+FPFH inputs. Each thread count runs all frontends except MAC$^\dagger$ sequentially in one job on one node, from the same code version. All rows use compiled finishing, with Open3D's point-to-plane ICP and Tukey weighting implemented in compiled loops and exact surface queries below $\tau$; the frozen finishing implementation takes 2.3\,s with \ls{} at four threads. Total includes the frontend, finishing and release moments; final selection takes additional milliseconds. Rows marked Compiled use compiled frontend kernels with growth restricted to each phase's initial support (\cref{S-sec:s-cost}). The dense switch in \cref{S-tab:s-cohortcost} does not engage on these pairs. The measured outcomes match the frozen implementation. SC$^2$-PCR and MAC use their authors' code, and TurboReg its author core. $^\dagger$The authors' C++ MAC fixes 12 OpenMP threads; it is timed only at 16 threads, in a later job on the same node, where rerun NGRR-LS and SC$^2$-PCR reproduce their times within 5\%. \main{\cref*{tab:main}} gives the 16-thread accuracy--time summary.}
\label{S-tab:s-cost}
\centering\footnotesize
\setlength{\tabcolsep}{4pt}
\ifdefined\supptablebox\else\newsavebox{\supptablebox}\fi
\sbox{\supptablebox}{\begin{tabular}{lrrrrrrr}
\toprule
 & \multicolumn{3}{c}{4 threads} & & \multicolumn{3}{c}{16 threads} \\
\cmidrule(lr){2-4}\cmidrule(lr){6-8}
Frontend & Frontend & Finishing & Total & & Frontend & Finishing & Total \\
\midrule
RANSAC-1M & 0.50 & 0.13 & 0.62 & & 0.44 & 0.10 & 0.54 \\
SC$^2$-PCR & 1.18 & 0.11 & 1.30 & & 0.83 & 0.09 & 0.92 \\
MAC$^\dagger$ & -- & -- & -- & & 1.55 & 0.11 & 1.68 \\
TurboReg & 2.72 & 0.17 & 2.91 & & 1.84 & 0.11 & 1.96 \\
\ngr{} & 4.03 & 0.13 & 4.17 & & 1.57 & 0.09 & 1.67 \\
\ls{} & 4.20 & 0.12 & 4.34 & & 1.73 & 0.09 & 1.82 \\
\midrule
\ngr{}, compiled & 0.65 & 0.13 & 0.78 & & 0.29 & 0.10 & 0.38 \\
\ls{}, compiled & 0.65 & 0.12 & 0.78 & & 0.30 & 0.09 & 0.39 \\
\bottomrule
\end{tabular}
}%
\ifdim\wd\supptablebox>\columnwidth\resizebox{\columnwidth}{!}{\usebox{\supptablebox}}\else\usebox{\supptablebox}\fi

\end{table}

\begin{table}[t]
\caption[Run time per cohort]{Median seconds per pair of the frozen and compiled \ls{} frontends and of the compiled finishing step on the first 100 pairs of each test cohort, at four threads; both implementations run one after the other in one job per cohort (jobs may share a node). Ratio: frozen over compiled frontend. The compiled frontend uses the frozen dense products where supports exceed 1600 matches or 70\% of the candidates are non-negligible (\cref{S-sec:s-cost}). These are separate runs from \cref{S-tab:s-cost}, so 3DLoMatch+FPFH differs slightly.}
\label{S-tab:s-cohortcost}
\centering\footnotesize
\setlength{\tabcolsep}{4pt}
\begin{tabular}{lrrrr}
\toprule
 & \multicolumn{3}{c}{\ls{} frontend} & \\
\cmidrule(lr){2-4}
Cohort & Frozen & Compiled & Ratio & Finishing \\
\midrule
3DMatch+FPFH & 3.74 & 0.65 & 5.7$\times$ & 0.11 \\
3DLoMatch+FPFH & 4.53 & 0.71 & 6.4$\times$ & 0.15 \\
KITTI+FPFH & 1.30 & 0.40 & 3.3$\times$ & 0.15 \\
ETH+FPFH & 4.59 & 0.63 & 7.2$\times$ & 0.19 \\
\midrule
3DMatch+FCGF & 4.27 & 2.32 & 1.8$\times$ & 0.09 \\
3DLoMatch+FCGF & 4.60 & 1.05 & 4.4$\times$ & 0.13 \\
KITTI+FCGF & 5.07 & 3.41 & 1.5$\times$ & 0.15 \\
3DMatch+GeoTr & 5.24 & 3.66 & 1.4$\times$ & 0.09 \\
3DLoMatch+GeoTr & 3.54 & 2.34 & 1.5$\times$ & 0.09 \\
\bottomrule
\end{tabular}


\end{table}

\paragraph{Locality and the compiled frontend.} \Cref{S-tab:s-cost} times every frontend at four and sixteen threads, each thread count in one job on one node. In the frozen \ls{} frontend, which is among the slowest implementations in this comparison, two dense products account for about 70\% of runtime: growth of 1024 phases and triplet scoring. Both products are local on 3DLoMatch+FPFH.

After growth, all but $10^{-6}$ of a phase's weight lies on a median of 100--190 out of up to 5000 matches (380--600 initially). A triplet score is non-negligible only for matches compatible with both members of its pair. Our compiled implementation grows each phase on its initial support: matches holding at least $10^{-6}$ of the largest initial weight. It gathers the block once for all phases of a seed pair. It forms each triplet score from the compatible matches only (terms below $10^{-9}$ of the pair's largest common weight and compatibilities below $10^{-10}$ are dropped) and scores candidates in the order of an upper bound, so that each batch's 3072 best are found without scoring the rest. The remaining steps run in compiled loops, and the libraries for the pair supports and the $3\times3$ SVDs are unchanged. This support-restricted implementation of \main{\cref*{sec:method}} runs on CPU. The \ls{} frontend then takes 0.65\,s at four threads and 0.30\,s at sixteen, against 1.18\,s and 0.83\,s for SC$^2$-PCR's released implementation.

\paragraph{Learned correspondences and the dense switch.} Learned correspondences have larger supports: on the FCGF and GeoTransformer cohorts of 3DMatch and 3DLoMatch, the initial supports cover a median of 24--99\% of the matches (ten pairs each). Gathering these blocks for each seed pair costs more than using dense products. The compiled frontend therefore keeps the frozen dense growth when a pair's initial supports average more than 1600 matches, and the frozen dense triplet scores when more than 70\% of the candidates of its first 64 shortlisted pairs are non-negligible. Both switch points were set on 30 calibration pairs per cohort, none of them among the timed pairs. On these 300 pairs, the switch takes 0.4\% more time than choosing the faster product for every pair. Moving the thresholds within 1500--1700 matches and 65--75\% adds at most 1.1\,s to the 415\,s. The switch never engages on 3DLoMatch+FPFH, KITTI+FPFH and ETH, grows densely on 1--5\% of the other FPFH pairs, and on 10--99\% of the learned-descriptor pairs. With it, the compiled frontend is 3.3--7.2 times faster than the frozen one on the four FPFH cohorts (0.40--0.71\,s at four threads) and 1.4--4.4 times faster on the five learned-descriptor cohorts (1.05--3.66\,s; \cref{S-tab:s-cohortcost}), where most of the time stays in the same dense products.

\paragraph{Same outcomes, different pools.} With compiled verification and the dense switch, \ls{} matches the frozen pipeline's outcomes on all 15\,035 pairs of the training cohort and the nine test cohorts (the same success or failure of the native output, the refined-pool ceiling, the finishing step's choice and the released candidate). The compiled \ngr{} also matches these outcomes on the 4781 pairs of the training cohort and 3DLoMatch+FPFH. Compiled triplet scoring selects different nuclei, so the candidate pools differ. Sparse-matrix kernels inside the dense products were slower than the dense products, since 7--17\% of the weights stay above any useful threshold. No two phases coincide, so duplicate merging saves no work.

\paragraph{Compiled verification.} All frontends share the same verification implementation. Surface evidence uses only nearest-neighbor distances below $\tau$, because larger distances contribute zero. A grid with that cell size therefore answers the queries exactly from 27 cells. ICP and greedy suppression run in compiled loops. The ICP reproduces Open3D's \texttt{registration\_icp} iterations: on 829 calls, every stage found the same number of correspondences and poses differed by at most $4\times10^{-6}$ degrees. Other nearest-neighbor queries use several threads, correspondence scoring runs in a compiled loop, and residual moments are computed in place. The resulting success/failure outcomes match the frozen implementation on every checked pair. Refined arrays differ by at most $10^{-12}$, the level of variation already seen between runs of the frozen implementation.

Finishing takes 0.12\,s instead of 2.3\,s with \ls{} at four threads (0.09--0.19\,s across test cohorts). On 3DLoMatch+FPFH, the timed pipeline takes 0.78\,s per pair at four threads and 0.39\,s at sixteen (SC$^2$-PCR with the same verification: 1.30\,s and 0.92\,s). These timings compare implementations, including the authors' C++ MAC; published results in \main{\cref*{tab:literature}} have unmatched time budgets.

\paragraph{Overheads of the units.} On the training cohort the site capacity adds 13\% to the frozen frontend time (mean time ratio 1.13), doubling the nuclei 64\%. Halving or doubling either capacity constant changes the refined-pool ceiling and the final count on 1000 training pairs by at most three pairs. With \ls{}, finishing uses on average 12.6 of its 32 ICP slots over the nine test cohorts; the pre-registered extension that fills the unused slots is not used (\cref{S-sec:s-protocol}). The release assay adds no ICP: one nearest-neighbor pass over the finishing step's query points accumulates the residual moments, and the rule itself takes milliseconds.

\paragraph{Actual use and validation scope.} The common refinement cap allows different actual ICP counts. On 3DLoMatch+FPFH, mean counts are approximately 21.9 for NGRR-LS, 22.7 for SC$^2$-PCR, 23.3 for MAC and 23.4 for TurboReg. Timings include actual work. The all-cohort checks compare success/failure indicators, without requiring numerically identical winning poses. Input loading and feature extraction are excluded from the registration timings; final release selection takes milliseconds beyond the in-place moment calculation.
